\section{Cross comparison in the slab}\label{sec:comparison}


The construction of Section~\ref{sec:opening} reduces the condensate
bound to the following estimate.

\begin{proposition}[Uniform second moment]\label{prop:second-moment}
After fixing the cutoffs as above, there is a constant $C_v<\infty$
such that, for all sufficiently small densities and sufficiently large
volumes at each density,
\[
                    \int f^2\dd\mu\le C_v.
\]
The constant is uniform for $\rho/2\le N/V\le2\rho$ and
$0<T\le\rho^2$.
\end{proposition}

We prove this by comparing the two changes in \eqref{eq:27} with a
single reverse insertion. Interactions away from encounters must cancel:
a fixed cost at every path cell would grow with the volume. The
analytic comparison below leaves only a fixed cost for shared labels
and a vanishing cost for nearby distinct bridges. To retain this gain
after averaging, we must also preserve the configuration-dependent
path and label choices. We establish that preservation before removing
their constraints from any conditional integral.

\subsection{The reverse insertion and the slot partition}

Fix a term of \eqref{eq:27}. The tuple $I'$ is selected along
$P'=P(C,\gamma')$ and produces the open state $O$; the density at $O$
uses a tuple $I$ along $P=P(\widehat C,\gamma)$. Their first label
$i_0$ is common. Write $e$ for the last label of $I$, and introduce
the reverse state $B$ by the following endpoint assignments:
\[
\begin{array}{c|c|c}
 \text{state}&\text{starting tuple}&\text{inner ending tuple}\\ \hline
 C&S^C&Y\\
 O&p_{I'}S^C&Y\text{ with }Y_{i_0}\text{ replaced by }z\\
 \widehat C&p_I^{-1}p_{I'}S^C&Y\\
 B&p_I^{-1}S^C&Y\text{ with }Y_e\text{ replaced by }x.
\end{array}
\]
The outer tuple $Y$ is unchanged in all four states. The table specifies
endpoints; the path variables retained or integrated in each state are
described next. Figure~\ref{fig:endpoints} shows the two single-tuple
assignments. In the cross term, the forward state uses $I'$ and the
reverse state uses $I$.
\begin{figure}[H]
\centering
\begin{tikzpicture}[x=1.35cm,y=1.15cm,>=Stealth,
 every node/.style={font=\small}]
 \begin{scope}
  \node at (1,1.55) {Forward insertion $O$};
  \foreach \j in {0,1,2} {
   \node (s\j) at (\j,1) {$S^C_{i_\j}$};
  }
  \node (y1) at (1,0) {$Y_{i_1}$};
  \node (y2) at (2,0) {$Y_{i_2}$};
  \node (z) at (2.7,0) {$z$};
  \node[gray] at (0,0) {$Y_{i_0}$};
  \draw[->] (s0)--node[midway,above right] {$i_1$} (y1);
  \draw[->] (s1)--node[midway,above right] {$i_2$} (y2);
  \draw[->,thick] (s2)--node[midway,right] {$i_0$} (z);
  \node at (0,-.5) {$a$};
  \node at (1,-.5) {middle cell};
  \node at (2.35,-.5) {$b$};
 \end{scope}
 \begin{scope}[xshift=6.0cm]
  \node at (1,1.55) {Reverse insertion $B$};
  \foreach \j in {0,1,2} {
   \node (t\j) at (\j,1) {$S^C_{i_\j}$};
  }
  \node (x) at (-.7,0) {$x$};
  \node (v0) at (0,0) {$Y_{i_0}$};
  \node (v1) at (1,0) {$Y_{i_1}$};
  \node[gray] at (2,0) {$Y_{i_2}$};
  \draw[->,thick] (t0)--node[midway,left] {$i_2$} (x);
  \draw[->] (t1)--node[midway,above left] {$i_0$} (v0);
  \draw[->] (t2)--node[midway,above left] {$i_1$} (v1);
  \node at (-.35,-.5) {$a$};
  \node at (1,-.5) {middle cell};
  \node at (2,-.5) {$b$};
 \end{scope}
\end{tikzpicture}
\caption{Endpoint assignments relative to the same original closed-state starts
$S^C$, for a tuple of three labels in successive
cells. The forward change opens slot $i_0$ at $z\in b$; the reverse
change opens slot $i_2$ at $x\in a$. Gray points remain outer endpoints
but cease to be inner endpoints. Each arrow is labeled by its current
slot and reassigns the original start at its tail to that slot. The two panels
illustrate the single-tuple changes separately; the second-moment comparison
may use different tuples for the forward and reverse changes. Arrows show
only assignments, not sampled
Brownian trajectories or a spatial straightness condition.}
\label{fig:endpoints}
\end{figure}


Identify a tuple with its slot set when taking intersections and
differences, and partition the selected slots as
\[
 \mathsf h=I\cap I',\quad
 \mathsf d=\mathsf h\cup\operatorname{pred}_I(\mathsf h),\quad
 \mathsf j=I\setminus\mathsf d,\quad
 \mathsf l=I'\setminus\mathsf h,\quad
 \mathsf L=(I\cup I')^c .
\]
Thus $\mathsf h$ contains the shared labels, while $\mathsf d$ also
contains their predecessors on $I$. In particular
\(|\mathsf d|\le2|\mathsf h|\) and $e\in\mathsf d$, because the
first label is shared. The two unshared portions that can cancel are
$\mathsf j$ and $\mathsf l$; $\mathsf L$ contains the untouched slots.
Let $J$ count cells of $P$ within distance $R_*$ of $P'$ and conversely,
adding both counts. Call these occurrences near and all others far.

All \(\mathsf L\) paths agree throughout. On \(\mathsf j\), \(C=O\) for endpoints and paths since the slots are off \(I'\), while \(B\) has the endpoints of \(\widehat C\): the \(I\)-successor is not shared, hence also off \(I'\), and \(S_i^B=S_{\operatorname{succ}_I(i)}^C=S_{\operatorname{succ}_I(i)}^O=S_i^{\widehat C}\) there. On \(\mathsf l\), \(B\) has the endpoints of \(C\), and \(\widehat C=O\). The separate path integrations off \(\mathsf L\) and our identification for the comparison are thus:
\[
\begin{array}{c|l|l}
 \text{set} & \text{path variables in (27)} & \text{use in }B\\ \hline
 \mathsf j & C_{\mathsf j}=O_{\mathsf j},\ \widehat C_{\mathsf j}
            & B_{\mathsf j}=\widehat C_{\mathsf j}\ \text{initially}\\
 \mathsf l & C_{\mathsf l},\ O_{\mathsf l}=\widehat C_{\mathsf l}
            & B_{\mathsf l}=C_{\mathsf l}\\
 \mathsf d & C_{\mathsf d},\ \widehat C_{\mathsf d},\ O_{\mathsf h}
            & B_{\mathsf d}\ \text{fresh}.
\end{array}
\]
We will also reintegrate some of \(\mathsf j\) below; \(O\) on \(\mathsf d\setminus\mathsf h\) equals \(C\) there.

We work first on the domain of the original nonzero weights and cuts in (27). All paths there on \(I\cup I'\) in \(C,O,\widehat C\) satisfy \(\mathrm{cut}_R\), by the two cuts in (27) and the two eligibilities. Displacements even for all endpoints needed in conditional integrations are at most \((D+10)s\): in \(O\) use the forward displacement rule from \(C\) on \(I'\) and from \(\widehat C\) on \(I\) (indeed \(S^O=p_I S^{\widehat C}\)). This also covers the non-selected slots on the union in \(C,\widehat C\), which individually have the same endpoints as in \(O\). For \(B\) on \(\mathsf d\), if the successor is shared use its start from an eligible slot of \(C\), close to the successor cell (for \(e\) use that \(x\) is in the first cell); otherwise the start agrees with \(\widehat C\).

For the beginning and end anchors of the integrated rows where (16) is needed the scores against \(\mathsf L\) are \(\le C_0\). In \(O\) the starts on \(I'\) come from eligible starts of \(C\), on \(I\) from eligible starts of \(\widehat C\); the respective end anchors are covered by eligibility or the dummy \(z\) test. Similarly \(B_{\mathsf j},B_{\mathsf l}\) use eligible endpoints from \(\widehat C,C\); on \(\mathsf d\) use the successor rule just given and end eligibility of \(\widehat C\) or the dummy \(x\) test. Paths in \(\mathsf L\) are the same in the scores. Also any pair between \(\mathsf j,\mathsf l\) in \(O\), and in \(B\), has starts separated by more than \(r\) and likewise ends: by the indicated eligible sources the starts used on both sets are isolated in the common start multiset, of which each of \(O,B\) assigns distinct elements (indices) to distinct slots. These ordinary end slots are isolated among \(Y\) by eligibility in \(\widehat C,C\) respectively (no separation from the exceptional inner end is needed).

All confined bridge measures use nearest lifts. Write \(\pi_i^X\) for the probability bridge on that lift before imposing the cut, and factor out the kernels \(k\) of those individual lifts (not the periodic sums), with product denoted \(K_X(G)\) on a set \(G\). Factor the kernels from the cut normalizations as well (write \(\overline{\mathcal Z}\) for the remaining normalization). The proposed-path kernels cancel, and relative to using baseline kernels the ratio left in (27) off \(\mathsf L\) is
\[
 \frac{K_C(I\cup I') K_{\widehat C}(I)}{K_B(I\cup I') K_O(I)}
      =\frac{K_C(\mathsf d)K_{\widehat C}(\mathsf d)}
             {K_B(\mathsf d)K_O(\mathsf d)}\le e^{C|\mathsf h|}.      \tag{28}\label{eq:28}
\]

\subsection{Conditional laws and the two denominator bounds}

For \(G=\mathsf j\) or \(\mathsf l\) write \(\Phi_G^X\) for the exponential of minus the interactions internal to \(G\) and with \(\mathsf L\) only, and \(Z[G,X]\) for its integral with the probability bridges and cuts \(R\), fixing \(\mathsf L\). Use empty-product conventions.

\paragraph{A one-path marginal bound.}
In these normalized group distributions an individual path marginal is bounded by \(C_4\pi_i^X\). The bound also holds if some other paths of the group are frozen on the cut, normalizing conditionally with their interactions. Indeed, conditional on all other paths, the individual normalization for its incident weight is at least \(p_*\) times the exponential of minus the expected interactions under its free bridge conditioned on the cut, by Jensen. Against \(\mathsf L\) use (16) divided by \(p_*\); other confined paths of the group can interact only from nearby end cells (distance at most \(2R+5\)), a bounded number of slots depending on \(R\) by simplicity, and the one-fixed-trajectory expectation bound applies (again divide by \(p_*\)). Thus this incident normalization is at least \(p_* e^{-C/p_*}\). Positivity \(v\ge0\) now gives a uniform conditional density upper bound \(C_4\) relative to \(\pi_i^X\), possibly large but independent of density after fixing the cut and score constants; average the bound for marginals. We use this observation for \(\mathsf j\) with the endpoints of \(C\) (equal to those of \(O\)) or \(B\), and \(\mathsf l\) with those of \(O\).

The needed lower normalizations are
\[
 \begin{split}
     \overline{\mathcal Z}_{I'}(O)&\ge
              Z[\mathsf l,O]\exp[-C|\mathsf h|-o_\rho(1)J],\\
     \overline{\mathcal Z}_{I}(O)&\ge
              Z[\mathsf j,C]\exp[-C|\mathsf h|-o_\rho(1)J].
 \end{split}                                                          \tag{29}\label{eq:29}
\]
Here the small coefficients tend to zero after all cut and eligibility
constants have been fixed. For the first, sample \(\mathsf l\) with its normalized group law, and \(\mathsf h\) independently with product free bridges conditioned on the cuts. Writing \(\mathcal V_{G,G'}\) for the summed interactions between disjoint groups (internal pairs once for \(\mathcal V_{G,G}\)), the ratio to \(Z[\mathsf l,O]\) is the product of free cut probabilities on \(\mathsf h\) times
\[
 \mathbb E\exp[-\mathcal V_{\mathsf l,\mathsf j}
               -\mathcal V_{\mathsf l,\mathsf d\setminus\mathsf h}
               -\mathcal V_{\mathsf h,\mathsf L\cup\mathsf j\cup(\mathsf d\setminus\mathsf h)\cup\mathsf l}
               -\mathcal V_{\mathsf h,\mathsf h}]
\]
with the complement frozen in \(O\). The interaction of \(\mathsf l\) with the frozen \(\mathsf j\) costs at most \(C/\sqrt r\) in expectation per nearby pair: it is zero in the boundary intervals by endpoint separation and both restraints (distance at least \(r-2r/10>R_v\)), and use marginal domination after those intervals. There are at most \(C(R)J\) possible pairs, since end cells must be within \(2R+5\) and are distinct within each set. Interaction with frozen \(\mathsf d\setminus\mathsf h\) costs at most \(C|\mathsf d|\), without needing the small bound. The independently sampled paths cost at most \(C|\mathsf h|\) with \(\mathsf L\) by (16) and with the other integrated or frozen paths by the fixed-trajectory bound and confinement (only a bounded number of slots in a fixed-radius end-cell neighborhood in these latter groups, allowing the special end). Jensen gives the first bound. 

For the second denominator, sample $\mathsf j$ with its group law
(the endpoints of $C$, which agree with those of $O$ on this set),
and sample $\mathsf d$ with independent free cut bridges. The frozen
selected complement is now $\mathsf l$. The missing interactions are
\[
 \mathcal V_{\mathsf j,\mathsf l},\qquad
 \mathcal V_{\mathsf d,\mathsf L\cup\mathsf j\cup\mathsf l},\qquad
 \mathcal V_{\mathsf d,\mathsf d}.
\]
The first sum costs $C r^{-1/2}J$, using the same endpoint separation
in $O$ and the marginal bound for $\mathsf j$. The other sums and the
free cut probabilities cost $C|\mathsf d|\le2C|\mathsf h|$. Jensen's
inequality gives the second line of \eqref{eq:29}.

\paragraph{The remaining normalized factors.}
Using \eqref{eq:28} and \eqref{eq:29}, drop interaction costs from the numerator weights by positivity except for: internal \(\mathsf L\); \(\Phi_{\mathsf j}^C,\Phi_{\mathsf l}^C\) from the closed weight; \(\Phi_{\mathsf l}^O\) from the proposal; \(\Phi_{\mathsf j}^{\widehat C}\) from the hypothetical old weight. 

To spell out the remaining integrations, let \(\pi_G^X\) denote probability-bridge products, \(c_G^X=\mathbf 1_{\mathrm{cut}_R(X_G)}\), and \(d\nu_G^X=c_G^X\Phi_G^X d\pi_G^X/Z[G,X]\) for the normalized groups. Besides the exponential comparison factors, the choice weights/constraints, \(K_B(I\cup I')\) and the original internal measure on \(\mathsf L\), the upper integrand on the nonzero domain now uses
\[
 (c_{\mathsf l}^C\Phi_{\mathsf l}^C\,d\pi_{\mathsf l}^C)
 (c_{\mathsf j}^{\widehat C}\Phi_{\mathsf j}^{\widehat C}\,d\pi_{\mathsf j}^{\widehat C})
 \ d\nu_{\mathsf j}^C\,d\nu_{\mathsf l}^O\,
 d\pi_{\mathsf d}^C\,d\pi_{\mathsf d}^{\widehat C}\,d\pi_{\mathsf h}^O .
\]


In particular the extras \(C_{\mathsf j},O_{\mathsf l}\) occur as normalized factors (fixing just endpoints and \(\mathsf L\) in their laws); all cuts inserted in these measures hold on the original domain. We still account for constraints and far choice weights below. For a near occurrence in cell \(Q\) bound the value of the choice factor on the nonzero domain by \([(1-\varepsilon)n_Q]^{-1}\) for count \(n_Q\) of that outer cell. We can and do carry the indicator of the original constraints (including both nonzero choices) even when bounding such values; in particular the counts along both paths are in the stated occupancy interval there.

\subsection{Restoring the baseline interaction weight}

We describe the weight restoration in \(B\). Put in \(\mathsf j_n\subseteq\mathsf j\) those with end cell within \(2R+5\) of some \(\mathsf l\) cell; other \(\mathsf j\) cannot interact with \(\mathsf l\) on the cuts. 

\paragraph{Fixed data and feasible fibers.}
Throughout this step, all endpoints, tuples, cell paths and skeletons
are fixed. The path integrations divide as follows:
\[
\begin{array}{c|l}
 \text{role}&\text{path variables}\\ \hline
 \text{fixed in both integrals}
   &\mathsf L,\quad B_{\mathsf l}=C_{\mathsf l},\quad
             B_{\mathsf j\setminus\mathsf j_n}
                       =\widehat C_{\mathsf j\setminus\mathsf j_n}\\
 \text{integrated in the original term}
   &\widehat C_{\mathsf j_n},\quad C_{\mathsf j},\quad O_{\mathsf l},
      \quad C_{\mathsf d},\quad\widehat C_{\mathsf d},\quad O_{\mathsf h}\\
 \text{integrated in the baseline term}&B_{\mathsf j_n},\quad B_{\mathsf d}.
\end{array}
\]
Call the fixed data \emph{feasible} if some choice of the remaining
original paths satisfies both nonzero selection weights, their specified
adaptations, and the cuts. Call such a choice a \emph{witness}. We only
need an upper comparison on feasible data: otherwise the original inner
integral vanishes. The probability laws of the extra original paths
depend only on the fixed endpoints and $\mathsf L$, not on
$\widehat C_{\mathsf j_n}$.

\paragraph{Restoring the missing interactions.}
First consider interaction weights alone. The bound needed for replacing
the $\widehat C_{\mathsf j_n}$ integral by a baseline integral is
\begin{equation}
 \Phi_{\mathsf l}^B
      \int c_{\mathsf j_n}^B\Phi_{\mathsf j}^B\,d\pi_{\mathsf j_n}^B
 \ \le\ e^{C|\mathsf h|+o_\rho(1)J}
      \int c_{\mathsf j_n}^B c_{\mathsf d}^B
       \exp\!\Big[-\sum_{\{i,i'\}\not\subseteq\mathsf L}
                          \mathcal V_{ii'}(B)\Big]\,
                       d\pi_{\mathsf j_n}^B d\pi_{\mathsf d}^B ,
\tag{29a}\label{eq:restoring-weights}
\end{equation}
where the sum is over unordered pairs of distinct slots indicated. 

Indeed, normalize
$c_{\mathsf j_n}^B\Phi_{\mathsf j}^B\,d\pi_{\mathsf j_n}^B$ with
$B_{\mathsf j\setminus\mathsf j_n}$ frozen, and independently sample
$\mathsf d$ by free bridges conditioned on their cuts. The earlier
one-path marginal bound applies also to this partially frozen group:
the score against $\mathsf L$ is still bounded, and only boundedly many
frozen or integrated $\mathsf j$ paths can be nearby. The ratio of the
full cut integral on the right to the product on the left is the free
cut probability product on $\mathsf d$, multiplied by the expectation
of the exponential of minus the missing interactions.

The $\mathsf j_n$--$\mathsf l$ interactions cost $o_\rho(1)J$ by
marginal domination, endpoint separation and the boundary restraints.
The frozen remainder of $\mathsf j$ cannot interact with $\mathsf l$
on the cuts. Interactions incident to $\mathsf d$ cost
$C|\mathsf d|$ by the free cut bounds, including \eqref{eq:16}.
Jensen and the cut probabilities prove \eqref{eq:restoring-weights}.
Its endpoint and frozen-path hypotheses follow from original feasibility,
independently of the reintegrated paths.

This inequality concerns interaction weights alone. To apply it to the
original integral, we must still account for the far selection factors
and the path-adaptation constraints. The next step proves that their
needed properties hold for every baseline cut replacement.

\subsection{Tracking the constraints in this comparison}

The original choices depend on which cells are usable and which slots
are eligible. We now compare those tests between any original witness
and every replacement of $B_{\mathsf j_n},B_{\mathsf d}$ on the cuts,
with the other data fixed as above. This universal comparison will
justify using \eqref{eq:restoring-weights} inside the constrained integral.

In $B$, use the outer tuple $Y$, including $Y_e$, for cell counts
and end isolation, and declare the exceptional slot $e$ ineligible.
Other eligibility tests use the starts and paths as before; the
exceptional path is included in all score sums. Let $S_0(B)$ be the
resulting unusable cells. Define $S_1(B)$ by the stricter score threshold
$C_0-\Delta_A$ and the stronger required eligible fraction
$1-\varepsilon/2$, with the same occupancy interval. Thus
$S_0(B)\subseteq S_1(B)$. Write $S(C),S(\widehat C)$ for the
unusable cells in the two witness states. The table compares the tests;
its last column gives the ineligible fraction causing failure when the
count lies in the allowed interval.
\[
\begin{array}{c|c|c|c}
 \text{bad-cell set}&\text{allowed count}&\text{score cutoff}
                                      &\text{failure fraction}\\ \hline
 S(C),S(\widehat C),S_0(B)&[b_-m,b_+m]&C_0&>\varepsilon\\
 S_1(B)&[b_-m,b_+m]&C_0-\Delta_A&>\varepsilon/2\\
 \text{defective cells}
    &[2b_-m,b_+m/2]&C_0-\Delta_A&\ge\varepsilon/4.
\end{array}
\]
Counts outside the displayed interval also cause failure. The symmetric
defective test uses actual inner endpoints, whereas $S_0(B),S_1(B)$
use the outer-end convention just specified. The witness states $C$
and $\widehat C$ are closed, so their inner ends are $Y$.

\begin{lemma}[Stability under every cut replacement]\label{lem:replacement-stability}
Fix feasible endpoint and path data as in the preceding subsection.
For every replacement of $B_{\mathsf j_n},B_{\mathsf d}$ on the cuts
and every original witness $C,\widehat C$, one has
\[
 \begin{gathered}
        S(C),S(\widehat C)\ \subseteq S_1(B),\\
        S(C)=S_0(B)\ \text{pointwise outside distance }R_*\text{ of }P,\\
        S(\widehat C)=S_0(B)\ \text{pointwise outside distance }R_*\text{ of }P'.
\end{gathered}                                                       \tag{30}\label{eq:30}
\]
In the far cells specified by the last two lines, all individual
eligibility indicators and eligible counts agree as well.
\end{lemma}
\begin{proof}
The state $C$ differs from $B$ only on $I$, and $\widehat C$ differs
from $B$ only on $\mathsf l,\mathsf d,\mathsf j_n$. All changed paths
in either comparison are confined. Their slots lie on at most two simple
cell paths, with the special inner end of $B$ in $a$.

At an unchanged slot's anchor, the change in either score is at most
$\Delta_A$. Indeed only changed paths with actual inner ends within
$(R+R_A)s$ of the anchor can contribute. In either state their number
is at most $2(2R+2R_A+4)^3+1$, by outer-cell counting and the one moved
exceptional end; also $\int_0^{\delta/2}w\le3$. Start isolation is
unchanged because the unordered start tuple is unchanged, and ordinary
outer-end isolation is unchanged as well. Consequently the eligible
count in either witness is at least the stricter eligible count of $B$
minus two per outer cell, allowing for the changed slots. If a cell
passes the $S_1(B)$ test, then its witness eligible count is at least
$(1-\varepsilon/2)n_Q-2\ge(1-\varepsilon)n_Q$ for sufficiently large
$m$. The occupancy counts are identical. This proves the inclusion
in \eqref{eq:30} by contraposition.

For exact equality, changed paths comparing with \(C\) lie near \(P\). Comparing with \(\widehat C\), their end cells lie within \(2R+5\) of \(P'\): note \(\mathsf d\) uses shared slots or their predecessors, including the last cell \(b\), and the special end is in \(a\). In a cell beyond the respective distance \(R_*\) there are no changed slots; an unchanged macro-conforming path has both anchors close enough to its cell that scores are exactly unchanged by the indicator cutoff in (15), and a macro failure stays ineligible anyway. The choice of \(R_*\) more than covers these distances. This proves equality of all actual eligibility flags and counts in the respective far cells.
\end{proof}

\paragraph{Lists containing every feasible pair.}
Define an explicit path predicate given \(B,\gamma,\gamma'\). Partition \(S_1(B)\) into components with adjacency distance up to \(K_b=2R_*+10\). Call components bridges if within \(R_*+2\) of both skeletons (set distances), and put \(b_{\rm tot}\) equal to the sum of their sizes. For the list of possible \(P\), take arbitrary subsets of cells in the bridge components as bad; from every nonbridge component intersecting \(\gamma\) take its intersection with \(S_0(B)\) as bad, ignoring the others. Run the deterministic cell-set adaptation for these sets, listing paths when it accepts geometrically. Similarly list \(P'\) using \(\gamma'\). Let \(G\) be the predicate that both given paths are on the respective lists.

The lists depend only on $B,\gamma,\gamma'$; $G$ tests membership of
the given cell-path pair and has no further dependence on its labels.

\begin{lemma}[Path lists and their encounter bound]\label{lem:path-lists}
For every original witness and every baseline cut replacement in
Lemma~\ref{lem:replacement-stability}, the witness paths $P,P'$ satisfy
$G$. For arbitrary $B$, the lists contain at most $4^{b_{\rm tot}}$
path pairs. Every listed pair satisfies
\[
                       J\le J_0+C(R_*)b_{\rm tot}.
                       \tag{31}\label{eq:32}
\]
Here $J_0$ counts encounters of the original skeletons, as in
Lemma~\ref{lem:skeleton}.
\end{lemma}
\begin{proof}
First fix a witness and any cut replacement. A nonbridge component of
$S_1(B)$ intersecting $\gamma$ stays outside distance $R_*$ of $P'$.
If it approached a skeleton vertex of $P'$, it would be a bridge. If it
approached a detour vertex, it would join, in the $K_b$ adjacency, the
potential component containing that detoured obstacle, which is hit by
$\gamma'$. It would again be a bridge. Therefore \eqref{eq:30} identifies
the actual bad subset of $\widehat C$ on this whole nonbridge component
as its intersection with $S_0(B)$.

On bridge components, choose the guessed subsets to be the actual bad
cells of $\widehat C$. Every actual nearest-neighbor bad component lies
in a single $K_b$-component of $S_1(B)$, and every such component hit by
$\gamma$ is retained. The reconstructed set therefore removes only whole
actual bad components not hit by $\gamma$. Lemma~\ref{lem:detour}
preserves acceptance and the adapted path under precisely these
deletions. Hence $P$ is listed. The same argument with $C,\gamma'$
gives $P'$. This reasoning holds for every replacement.

There are at most $2^{b_{\rm tot}}$ guesses for each path, giving
$4^{b_{\rm tot}}$ pairs. To bound their encounters, first count near
occurrences that lie on their own skeleton and have an opposite-skeleton
vertex within distance $R_*$. Since the adapted paths are simple, these
occurrences inject into the corresponding counts of $J_0$.

Every remaining near occurrence is a detour vertex or is close to an
opposite detour vertex. Any detour participating in such an encounter
belongs to a bridge component: it is adjacent to its obstacle, which
is hit by its skeleton, and either lies close to the other skeleton or
joins the component of an opposite detoured obstacle in the $K_b$
adjacency. Across both paths, these detours use at most $12b_{\rm tot}$
vertices by Lemma~\ref{lem:detour}. Their fixed-radius neighborhoods
contain at most $C(R_*)b_{\rm tot}$ cells. Simplicity bounds the number
of path occurrences in those cells by twice this number. This proves
\eqref{eq:32} with coefficient one on $J_0$. The argument also applies
to guessed subsets: their nearest-neighbor obstacles remain within
components of $S_1(B)$.
\end{proof}

\paragraph{The common far-choice product.}
For a slot, its choice factor in $B$ is its eligible indicator divided
by the eligible count in its outer cell, using the $S_0(B)$ criteria
and assigning zero when the count is zero. Let $H_F(B)$ be the product
of these factors for the selected far occurrences of both paths.
Unlike the path lists, this product depends on the selected labels.
Its indicators and denominators are nevertheless all evaluated on the
same fixed baseline configuration.

Lemma~\ref{lem:replacement-stability} identifies these factors exactly
with the original far factors: $P$ uses $\widehat C$ and $P'$ uses $C$.
At feasible fixed data their product therefore has a common value
$H_*$ for all original witnesses and all cut replacements. Indeed,
the equality holds between \emph{any} witness and \emph{every}
replacement, and the replacement cuts are nonempty. This is the
selection property needed to apply \eqref{eq:restoring-weights}.



Explicitly, at those fixed variables let \(\mathcal D\) impose the original remaining constraints, \(H_{\rm old}\) be the far product there, and use the product probability measure
\[
d\mathcal P=d\nu_{\mathsf j}^C\, d\nu_{\mathsf l}^O\,
 d\pi_{\mathsf d}^C\,d\pi_{\mathsf d}^{\widehat C}\,d\pi_{\mathsf h}^O
\]
on the extras (its laws do not depend on the integrated \(\widehat C_{\mathsf j_n}\)). Having bounded the near factors, the needed inequalities on feasibility, omitting the cut indicators equal to one on the fixed paths in \(\mathsf j,\mathsf l\), are
\[
\begin{split}
 &\Phi_{\mathsf l}^B
 \int c_{\mathsf j_n}^{\widehat C}\Phi_{\mathsf j}^{\widehat C}
       \mathbf 1_{\mathcal D} H_{\rm old}\,d\pi_{\mathsf j_n}^{\widehat C}\,d\mathcal P\\
 &\quad\le H_*\,\Phi_{\mathsf l}^B
      \int c_{\mathsf j_n}^B\Phi_{\mathsf j}^B\,d\pi_{\mathsf j_n}^B\\
 &\quad\le e^{C|\mathsf h|+o_\rho(1)J}
      \int c_{\mathsf j_n}^B c_{\mathsf d}^B
       \exp\!\Big[-\sum_{\{i,i'\}\not\subseteq\mathsf L}
                          \mathcal V_{ii'}(B)\Big]\mathbf 1_G H_F(B)\,
                       d\pi_{\mathsf j_n}^B d\pi_{\mathsf d}^B .
\end{split}
\]
The first inequality removes the original constraints only after
factoring out their common far product. The second uses
\eqref{eq:restoring-weights} and Lemma~\ref{lem:path-lists}:
$\mathbf1_GH_F(B)=H_*$ everywhere on the cut replacements. At
infeasible data the original inner integral is zero, so its comparison
with the nonnegative baseline integral is automatic. Thus the argument
compares inner integrals at fixed data, without choosing replacements
according to the number of witnesses they admit.

For precision, denote the original conditional path integral by
\[
\begin{split}
 \mathcal I_{I,I'}={}&\int W(C)h_{\gamma'}(C,I')
 \frac{\mathbf1_{\mathrm{cut}_R(O_{I'})}W_{I'}^O}
      {\mathcal Z_{I'}(O)}\\
 &\hspace{5mm}\cdot
 \frac{\mathbf1_{\mathrm{cut}_R(O_I)}}{\mathcal Z_I(O)}
                    W_I^{\widehat C}h_\gamma(\widehat C,I).
\end{split}
\]
The integrations are over the original closed paths, proposed paths on
$I'$, and hypothetical old paths on $I$, with the endpoints, tuples, and
skeletons fixed. Let $\mathbf1_{\mathrm{occ}}$ require that the outer
counts lie in $[b_-m,b_+m]$ in every cell of both paths.

\begin{lemma}[Integrated cross comparison]\label{lem:cross}
Fix two simple star cell paths from $a$ to $b$, their label tuples
matching the outer $Y$-cells and having a common first label,
two skeletons, and endpoint data $Y,S^C,x,z$ as in \eqref{eq:27}.
With $B$, the near occurrences, $G$, and $H_F(B)$ defined above, one has
\[
\begin{split}
 \mathcal I_{I,I'}\le{}&
 e^{C|\mathsf h|+\sigma_\rho J}\mathbf1_{\mathrm{occ}}
 \prod_{\substack{\text{near occurrences}\\\text{in cell }Q}}
                              [(1-\varepsilon)n_Q(Y)]^{-1}\\
 &\hspace{8mm}\cdot\int W(S^B,E^B;d\omega)\,\mathbf1_G H_F(B).
\end{split}\tag{32}\label{eq:31}
\]
Here $C<\infty$ is fixed and $\sigma_\rho\ge0$ tends to zero as
$\rho\downarrow0$ after the cutoffs are fixed, uniformly in all permitted
volumes and temperatures. A cell on both paths is counted twice in the
product. The right side is zero when the occupancy condition fails.
\end{lemma}
\begin{proof}
The preceding inner-integral comparison proves the displayed bound on
every feasible fiber. On an infeasible fiber the original contribution
is zero. We explain the remaining bookkeeping to specify precisely how
the full integral is recovered.
Keep $\mathsf L,B_{\mathsf l},B_{\mathsf j\setminus\mathsf j_n}$ fixed
and integrate the extensions in $\mathsf j_n,\mathsf d$ as in the
displayed ledger. Whenever the original domain is nonempty, all the
nearest-kernel factorizations apply. Restore the baseline kernels from
\eqref{eq:28}, then integrate the fixed paths. Finally drop the remaining
cut and nearest-lift restrictions from the baseline integral by
positivity, obtaining \eqref{eq:31}. Its predicates are measurable
functions of the baseline paths, specified by finite cell algorithms.
This use of the conditional comparison is valid pointwise at the fixed
data: a positive original inner integral has a witness, and otherwise
the required upper bound is automatic.

\end{proof}

The conditional comparison is now established. To average its path-list
cost, we need the defect bound under the reverse insertion. Under
$\mu_{ab}$, with $e$ the opened slot in each term, every specified
cell set $\mathcal U$ satisfies
\[
                  \mu_{ab}(\mathcal U\subseteq S_1(B))
                              \le C M^{-2} q^{|\mathcal U|},
                  \qquad q=q_{\rm cl}^{1/4},                          \tag{33}\label{eq:33}
\]
including the empty set bound. To compare to the symmetric defective test replace the exceptional outer endpoint by its actual inner end for counting and isolation. Counts change by at most one; only a bounded number of other end-isolation flags can change: those isolated before but not after are mutually separated by more than \(r\) and within \(r\) of the new point, and analogously in reverse (use packing). Their other tests are unchanged. Allow also the exceptional slot itself. For a potential bad cell with count \(n_Q(Y)\) in the occupancy interval, the ineligible count in the actual stricter-score test is thus at least \((\varepsilon/2)n_Q(Y)-C\), giving at least fraction \(\varepsilon/4\) for large \(m\); counts outside the interval imply the widened count defect after any change by one. This works pointwise in all the specified cells. Now apply \eqref{eq:22} and \eqref{eq:23}.

\subsection{Summing the comparison}

The conditional comparison is now complete, including its selection
predicates. We sum first over labels at fixed baseline paths, then over
the finite lists of possible cell paths. The remaining environment
average will be controlled by the insertion defect bound and the two
skeleton moments.

\paragraph{Summing label choices.}
Sum \eqref{eq:31} over candidate cell-path pairs and their tuples; this
candidate range depends on $a,b,Y$, but not on the starts. For each
tuple pair change variables from $S^C$ to $S^B$. The outer kernel is
invariant. Group terms by the last label $e$ of $I$, whose outer end
lies in $b$. At fixed $e$ and baseline configuration, the baseline
measure, $S_0,S_1$, and the path lists are independent of every other
label choice.

Fix a listed pair satisfying the occupancy bound. A far cell occurs on
only one path, and summing its choice factor from $H_F(B)$ costs at
most one. Each near occurrence contributes a factor
$(1-\varepsilon)^{-1}$. After removing these factors, the remaining
cellwise sums are
\[
\begin{array}{c|c|c}
 \text{cell}&\text{label choices}&\text{upper bound}\\ \hline
 \text{ordinary common cell }Q&\text{two independent labels}
       &1+(e^C-1)/n_Q\\
 \text{first cell }a&\text{one shared label}&e^C/n_a\\
 \text{last cell }b&\text{first tuple's label fixed to }e
       &[1+(e^C-1)/n_b]/n_b.
\end{array}
\]
For example, at an ordinary common cell, the $n_Q$ coincident choices
carry $e^C$ and the $n_Q(n_Q-1)$ distinct choices carry one, all divided
by $n_Q^2$. At $a$ the two tuples must share their first label; at $b$
one of the two labels is fixed by the grouping. The endpoint cells are
distinct. Simplicity makes these sums independent across cells, since
labels in different outer cells cannot coincide. Using $n_Q\ge b_-m$
and including $e^{\sigma_\rho J}$ gives the bound
\[
                           \frac{C}{m^2}e^{\kappa_0 J},
                       \qquad \kappa_0=2\varepsilon+o_\rho(1).
\]
Lemma~\ref{lem:path-lists} now sums the cell-path pairs, giving
$C m^{-2}\exp(\kappa_0J_0+C'b_{\rm tot})$ with fixed $C'$.

\paragraph{Recovering the reverse insertion.}
There is no remaining dependence on the dummy-for-\(B\) point \(z\) within \(b\); its base integral there gives \(1/M\). Together with the prefactor of (27) we have \(m/(m^2 M)=1/N\); summing \(e\) with \(Y_e\) in \(b\) and integrating the baseline with \(x\) in \(a\) using \(dx/V\) now gives exactly $\mu_{ab}$, with outer endpoint in $b$ and new
inner endpoint in $a$, as defined in Section~\ref{sec:defects}. Thus
\[
                 \int f^2\dd\mu
         \le C\sum_{a,b\ {\rm valid}} \mathbb E_{\gamma,\gamma'}\!
               \left[e^{\kappa_0 J_0}
                   \int e^{C' b_{\rm tot}}\,d\mu_{ab}\right].         \tag{34}\label{eq:34}
\]

\paragraph{The exponential moment of rare components.}
For fixed skeletons, we claim
\[
 \int e^{C'b_{\rm tot}}\,d\mu_{ab}
 \le \frac C{M^2}\exp\!\left[\vartheta(q)S_{K_b}\right],
 \qquad \vartheta(q)\longrightarrow0\quad(q\downarrow0),
\]
where $S_{K_b}$ is the soft-distance sum of Lemma~\ref{lem:skeleton}.
This is the estimate that allows the environment average in
\eqref{eq:34} to use the second skeleton moment.

Expand the product over realized bridge components of $S_1(B)$ with
factors $1+(e^{C'|T|}-1)$. Bound it by the sum over unordered
collections of disjoint connected sets $T$ in the $K_b$ adjacency,
each within $R_*+2$ of both skeletons, requiring their cells to lie in
$S_1(B)$. Apply \eqref{eq:33} to each union. Summing the products and
then dropping disjointness gives
\[
                   C M^{-2}
                     \exp\Big(\sum_T(e^{C'}q)^{|T|}\Big)
\]
(the sum in the exponent is over individual qualifying sets). 

Root their counting near vertex occurrences \(z\in\gamma\); there are a bounded number of root cells per occurrence and at most \(C(K_b)^j\) connected sets of size \(j\) per root by the bounded degree. To reach the other skeleton requires \(d_\infty(z,\gamma')\le 2(R_*+2)+K_b j\). Consequently the exponent sum is at most
\[
                        o_q(1)\sum_{z\in\gamma}
                                     e^{-d_\infty(z,\gamma')/K_b}
\]
by a geometric series: per root sum \((C(K_b)e^{C'}q)^j\) only at sizes with \(d_\infty(z,\gamma')/K_b\le j+2(R_*+2)/K_b\), extracting the displayed distance decay. 

The geometric-series bound proves the claim with a coefficient
$\vartheta(q)\to0$. Take density small enough that $2\kappa_0\le\alpha$
and $2\vartheta(q)\le\xi(K_b)$. Cauchy--Schwarz and
\eqref{eq:24}--\eqref{eq:25} then give, for each valid pair,
\[
 \mathbb E_{\gamma,\gamma'}
 \left[e^{\kappa_0J_0}\int e^{C'b_{\rm tot}}\,d\mu_{ab}\right]
 \le \frac C{M^2}
       \bigl(\mathbb Ee^{2\kappa_0J_0}\bigr)^{1/2}
       \bigl(\mathbb Ee^{2\vartheta(q)S_{K_b}}\bigr)^{1/2}
 \le \frac C{M^2}.
\]
There are at most $M^2$ ordered pairs. Substitution into
\eqref{eq:34} proves Proposition~\ref{prop:second-moment}.
